
\documentclass[11pt,a4paper]{article}

\usepackage[margin=0.85in]{geometry}
\usepackage{amsmath,amssymb}
\usepackage{enumitem}

\setlength{\parindent}{0pt}
\setlength{\parskip}{6pt}

\begin{document}

\section*{Problem}

There are two states, Red and Blue, initially with 100
residents each. Every resident supports exactly one of
two parties, Democrat or Republican. Initially, 60
residents of Blue and 30 residents of Red are Democrats.

Ten residents move from Blue to Red, and nobody
changes parties. Let $k$ be the number of Democrats
among the ten people who move.

Whenever a resident is selected, the selection is
uniform over all 200 people. Let $D$ be the event
that the selected resident is a Democrat, and let
$B$ be the event that the resident lives in Blue
at the time of selection.

\begin{enumerate}[label=(\alph*)]

\item Find $P(D\mid B)$, $P(D\mid B^c)$, and
$P(D)$ before the move. Does the move change
$P(D)$? Explain.

\item Determine all integer values of $k$ for
which the proportion of Democrats increases
strictly in both states.

Explain, using the law of total probability with
the appropriate weights, how both proportions
can increase while the overall proportion
stays unchanged.

\item For this part, take $k=5$ and consider
the population after the move. A survey
interviews the selected resident only if their
name appears on a list containing all Blue
residents and all Democrats in Red.

Let $A$ be the event that the selected resident
is on this list. Find $P(D\mid A)$ using the
law of total probability conditional on $A$,
with the partition $\{B,B^c\}$, and state
the two group weights.

Explain why this probability differs from the
proportion of Democrats in the whole population.

\end{enumerate}

\section*{Solution}

\subsection*{(a)}

Before the move,

\[
\because\quad
P(B)=P(B^c)=\frac{100}{200}=\frac12
\]

\[
P(D\cap B)=\frac{60}{200},
\qquad
P(D\cap B^c)=\frac{30}{200}.
\]

By conditional probability,

\[
\begin{aligned}
P(D\mid B)
&=\frac{P(D\cap B)}{P(B)}\\
&=\frac{60/200}{100/200}\\
&=\frac35
\end{aligned}
\]

Similarly,

\[
\begin{aligned}
P(D\mid B^c)
&=\frac{P(D\cap B^c)}{P(B^c)}\\
&=\frac{30/200}{100/200}\\
&=\frac3{10}
\end{aligned}
\]

By the law of total probability,

\[
\begin{aligned}
P(D)
&=P(D\mid B)P(B)
 +P(D\mid B^c)P(B^c)\\
&=\frac35\cdot\frac12
 +\frac3{10}\cdot\frac12\\
&=\frac9{20}
\end{aligned}
\]

After the move,

\[
\begin{aligned}
P'(B)&=\frac{90}{200},\\
P'(B^c)&=\frac{110}{200},\\
P'(D\mid B)&=\frac{60-k}{90},\\
P'(D\mid B^c)&=\frac{30+k}{110}.
\end{aligned}
\]

\[
\begin{aligned}
\therefore\quad P'(D)
&=P'(D\mid B)P'(B)
 +P'(D\mid B^c)P'(B^c)\\
&=\frac{60-k}{90}\cdot\frac{90}{200}
 +\frac{30+k}{110}\cdot\frac{110}{200}\\
&=\frac{60-k+30+k}{200}\\
&=\frac{90}{200}\\
&=\frac9{20}.
\end{aligned}
\]

\[
\Longrightarrow\quad
P'(D)=P(D)=\frac9{20}
\]

The move changes neither the total population
nor the total number of Democrats.

\subsection*{(b)}

After the move, the Democratic proportions are

\[
P'(D\mid B)=\frac{60-k}{90},
\qquad
P'(D\mid B^c)=\frac{30+k}{110}.
\]

For Blue,

\[
\begin{aligned}
\frac{60-k}{90}&>\frac{60}{100}\\
60-k&>54\\
k&<6.
\end{aligned}
\]

For Red,

\[
\begin{aligned}
\frac{30+k}{110}&>\frac{30}{100}\\
30+k&>33\\
k&>3.
\end{aligned}
\]

\[
\because\quad
3<k<6,\qquad k\in\mathbb Z
\]

\[
\Longrightarrow\quad
k\in\{4,5\}
\]

Before the move, the population weights are

\[
P(B)=\frac12,
\qquad
P(B^c)=\frac12.
\]

\[
\begin{aligned}
\therefore\quad P(D)
&=\frac35\cdot\frac12
 +\frac3{10}\cdot\frac12\\
&=\frac9{20}.
\end{aligned}
\]

After the move, the population weights become

\[
P'(B)=\frac9{20},
\qquad
P'(B^c)=\frac{11}{20}.
\]

By the law of total probability,

\[
\begin{aligned}
P'(D)
&=P'(D\mid B)P'(B)
 +P'(D\mid B^c)P'(B^c)\\
&=\frac{60-k}{90}\cdot\frac9{20}
 +\frac{30+k}{110}\cdot\frac{11}{20}\\
&=\frac{60-k}{200}
 +\frac{30+k}{200}\\
&=\frac9{20}.
\end{aligned}
\]

For example, when $k=5$,

\[
\begin{aligned}
P'(D\mid B)
&=\frac{55}{90}
>\frac{60}{100},\\[4pt]
P'(D\mid B^c)
&=\frac{35}{110}
>\frac{30}{100}.
\end{aligned}
\]

However,

\[
\begin{aligned}
P'(D)
&=\frac{55}{90}\cdot\frac{90}{200}
 +\frac{35}{110}\cdot\frac{110}{200}\\
&=\frac{55+35}{200}\\
&=\frac9{20}.
\end{aligned}
\]

\[
\Longrightarrow\quad
P'(D)=P(D)=0.45
\]

Both conditional proportions increase, but
their population weights change from
$(0.5,0.5)$ to $(0.45,0.55)$.

The larger weight is assigned to the state
with the lower Democratic proportion.
The change in weights offsets the increases
in the conditional proportions.

\subsection*{(c)}

\[
\because\quad k=5
\]

After the move,

\[
\begin{array}{c|cc}
&\text{Blue}&\text{Red}\\ \hline
\text{Population}&90&110\\
\text{Democrats}&55&35\\
\text{Republicans}&35&75
\end{array}
\]

The list contains all 90 Blue residents
and the 35 Democrats in Red.

\[
\therefore\quad |A|=90+35=125.
\]

The two group weights conditional on $A$ are

\[
\begin{aligned}
P(B\mid A)
&=\frac{90}{125}
=\frac{18}{25}\\[6pt]
P(B^c\mid A)
&=\frac{35}{125}
=\frac7{25}
\end{aligned}
\]

For the Blue group,

\[
\because\quad
A\cap B=B
\]

\[
\therefore\quad
P(D\mid A\cap B)
=\frac{55}{90}
=\frac{11}{18}.
\]

For the Red group, every resident on the
list is a Democrat.

\[
\therefore\quad
P(D\mid A\cap B^c)=1.
\]

By the conditional law of total probability,

\[
\begin{aligned}
P(D\mid A)
&=P(D\mid A\cap B)P(B\mid A)\\
&\quad+
P(D\mid A\cap B^c)P(B^c\mid A)\\[4pt]
&=\frac{55}{90}\cdot\frac{90}{125}
 +1\cdot\frac{35}{125}\\[4pt]
&=\frac{55+35}{125}\\
&=\frac{90}{125}.
\end{aligned}
\]

\[
\Longrightarrow\quad
P(D\mid A)=\frac{18}{25}=0.72
\]

For the whole population,

\[
P(D)=\frac{90}{200}=0.45.
\]

\[
\because\quad
|D\cap A|=90,\qquad
|A|=125,\qquad
|\Omega|=200
\]

\[
\therefore\quad
P(D\mid A)=\frac{90}{125}
>
\frac{90}{200}=P(D).
\]

The list excludes all 75 Republicans in Red
but includes every Democrat.

Thus, the selected group overrepresents
Democrats relative to the whole population.

\[
\Longrightarrow\quad
P(D\mid A)=0.72>0.45=P(D)
\]


\section*{Problem 2}

A sequence of $n\geq 1$ mutually independent trials is
performed. Each trial results in either success or failure,
and the probability of success on trial $i$ is $p_i$,
where $0\leq p_i\leq 1$.
Let $A_n$ be the event that the number of successes
among the first $n$ trials is even.

\textbf{(a)}
Find $P(A_2)$ by listing the possible outcomes,
and show that
\[
P(A_2)=
\frac{1+(1-2p_1)(1-2p_2)}{2}.
\]

\textbf{(b)}
By considering whether the last trial succeeds or fails,
obtain a recursion for $P(A_n)$ for $n\geq 2$.
Use the recursion to prove that
\[
P(A_n)=
\frac{1+\prod_{i=1}^{n}(1-2p_i)}{2}.
\]

You may use mathematical induction or repeated substitution
in the recursion. State the initial value explicitly.

\textbf{(c)}
Under the independence assumption, is
$P(A_n)=1/2$ if and only if at least one $p_i=1/2$?
Justify your answer.
Then remove the independence assumption and specify
probabilities for the four possible outcomes of two trials
so that $p_1=p_2=1/2$ but $P(A_2)\neq 1/2$.


\subsection*{Solution}

\textbf{(a)}

Let $S$ and $F$ denote success and failure,
respectively.

The four possible outcomes are
\[
\{SS,SF,FS,FF\}.
\]

\[
\begin{aligned}
P(SS)&=p_1p_2,\\
P(SF)&=p_1(1-p_2),\\
P(FS)&=(1-p_1)p_2,\\
P(FF)&=(1-p_1)(1-p_2).
\end{aligned}
\]

\[
A_2=\{SS,FF\}.
\]

\[
\begin{aligned}
\therefore\quad P(A_2)
&=P(SS)+P(FF)\\
&=p_1p_2+(1-p_1)(1-p_2)\\
&=1-p_1-p_2+2p_1p_2\\
&=\frac{2-2p_1-2p_2+4p_1p_2}{2}\\
&=
\frac{1+(1-2p_1)(1-2p_2)}{2}
\end{aligned}
\]



\textbf{(b)}

Let $S_n$ and $F_n$ denote success and failure
on the $n$-th trial.

\[
A_n=
(A_{n-1}\cap F_n)
\cup
(A_{n-1}^{c}\cap S_n).
\]

\[
\begin{aligned}
\therefore\quad P(A_n)
&=P(A_{n-1}\cap F_n)
 +P(A_{n-1}^{c}\cap S_n)\\
&=P(A_{n-1})(1-p_n)
 +[1-P(A_{n-1})]p_n\\
&=(1-2p_n)P(A_{n-1})+p_n.
\end{aligned}
\]

\textbf{Initial value:}

\[
P(A_1)=1-p_1.
\]

\textbf{Repeated substitution:}

\[
\begin{aligned}
P(A_n)-\frac12
&=(1-2p_n)P(A_{n-1})+p_n-\frac12\\
&=(1-2p_n)\left(P(A_{n-1})-\frac12\right).
\end{aligned}
\]

\[
\begin{aligned}
\therefore\quad P(A_n)-\frac12
&=(1-2p_n)\left(P(A_{n-1})-\frac12\right)\\
&=(1-2p_n)(1-2p_{n-1})
  \left(P(A_{n-2})-\frac12\right)\\
&=\cdots\\
&=\prod_{i=2}^{n}(1-2p_i)
  \left(P(A_1)-\frac12\right)\\
&=\prod_{i=2}^{n}(1-2p_i)
  \left(\frac12-p_1\right)\\
&=\frac12\prod_{i=2}^{n}(1-2p_i)(1-2p_1)\\
&=\frac12\prod_{i=1}^{n}(1-2p_i).
\end{aligned}
\]

\[
\Rightarrow\quad
P(A_n)=
\frac{1+\prod_{i=1}^{n}(1-2p_i)}{2}
\]


\textbf{(c)}

\textbf{Under independence:}

\[
P(A_n)=
\frac{1+\prod_{i=1}^{n}(1-2p_i)}{2}.
\]

\[
\begin{aligned}
P(A_n)=\frac12
&\Longleftrightarrow
1+\prod_{i=1}^{n}(1-2p_i)=1\\
&\Longleftrightarrow
\prod_{i=1}^{n}(1-2p_i)=0\\
&\Longleftrightarrow
\exists\,i\in\{1,\ldots,n\},
\quad 1-2p_i=0\\
&\Longleftrightarrow
\exists\,i\in\{1,\ldots,n\},
\quad p_i=\frac12.
\end{aligned}
\]

\[
\therefore\quad
P(A_n)=\frac12
\Longleftrightarrow
\exists\,i,\quad p_i=\frac12
\]

\textbf{Without independence:}

Consider the joint distribution
\[
\begin{aligned}
P(SS)&=\frac12,\\
P(SF)&=0,\\
P(FS)&=0,\\
P(FF)&=\frac12.
\end{aligned}
\]

\[
\begin{aligned}
p_1
&=P(SS)+P(SF)\\
&=\frac12+0=\frac12,
\end{aligned}
\]

\[
\begin{aligned}
p_2
&=P(SS)+P(FS)\\
&=\frac12+0=\frac12.
\end{aligned}
\]

\[
\begin{aligned}
P(A_2)
&=P(SS)+P(FF)\\
&=\frac12+\frac12\\
&=1\neq\frac12.
\end{aligned}
\]

\[
\because\quad
P(SS)=\frac12
\neq p_1p_2=\frac14,
\]

\[
\therefore\quad
\text{the two trials are not independent}.
\]

\[
p_1=p_2=\frac12,
\qquad
P(A_2)=1\neq\frac12.
\]



\section*{Problem 3}

To battle against spam, Bob installs two anti-spam programs.
An email is either legitimate (event $L$) or spam (event $L^c$).
Let $M_j$ be the event that program $j$ marks the email as
legitimate, for $j=1,2$. Assume that
\[
P(L)=0.10,\qquad
P(M_j\mid L)=P(M_j^c\mid L^c)=0.90
\quad (j=1,2).
\]

The two programs process the same email.
Their outputs are conditionally independent given $L$,
and also conditionally independent given $L^c$.

\begin{enumerate}[label=(\alph*)]

\item Find $P(L\mid M_1)$ and simplify your answer.

\item Find $P(L\mid M_1\cap M_2)$ and
$P(L\mid M_1\cap M_2^c)$.
Compare each with the prior probability $P(L)$.

\item Are $M_1$ and $M_2$ independent without
conditioning on the email's true status?
Prove your answer by computing the required probabilities.
Explain why your answer is consistent with the stated
conditional independence.

\end{enumerate}

\section*{Solution}

\noindent
\textbf{(a)}

\[
\because\quad
P(L^c)=0.90,\qquad
P(M_1\mid L^c)=0.10.
\]

By the law of total probability,
\begin{align*}
P(M_1)
&=P(M_1\mid L)P(L)
 +P(M_1\mid L^c)P(L^c)\\
&=0.90(0.10)+0.10(0.90)\\
&=0.18.
\end{align*}

By Bayes' theorem,
\begin{align*}
P(L\mid M_1)
&=\frac{P(M_1\mid L)P(L)}{P(M_1)}\\
&=\frac{0.90(0.10)}{0.18}\\
&=\frac{1}{2}.
\end{align*}

\[
\therefore\quad P(L\mid M_1)=0.50.
\]


\bigskip
\noindent
\textbf{(b)}

\textbf{Case 1:} $P(L\mid M_1\cap M_2)$.

By conditional independence,
\begin{align*}
P(M_1\cap M_2\mid L)
&=P(M_1\mid L)P(M_2\mid L)\\
&=0.90^2=0.81,\\[4pt]
P(M_1\cap M_2\mid L^c)
&=P(M_1\mid L^c)P(M_2\mid L^c)\\
&=0.10^2=0.01.
\end{align*}

\begin{align*}
P(M_1\cap M_2)
&=P(M_1\cap M_2\mid L)P(L)\\
&\quad+P(M_1\cap M_2\mid L^c)P(L^c)\\
&=0.81(0.10)+0.01(0.90)\\
&=0.09.
\end{align*}

By Bayes' theorem,
\begin{align*}
P(L\mid M_1\cap M_2)
&=\frac{P(M_1\cap M_2\mid L)P(L)}
{P(M_1\cap M_2)}\\
&=\frac{0.81(0.10)}{0.09}\\
&=0.90.
\end{align*}

\[
\therefore\quad
P(L\mid M_1\cap M_2)=0.90>P(L)=0.10.
\]


\medskip
\textbf{Case 2:} $P(L\mid M_1\cap M_2^c)$.

By conditional independence,
\begin{align*}
P(M_1\cap M_2^c\mid L)
&=P(M_1\mid L)P(M_2^c\mid L)\\
&=0.90(0.10)=0.09,\\[4pt]
P(M_1\cap M_2^c\mid L^c)
&=P(M_1\mid L^c)P(M_2^c\mid L^c)\\
&=0.10(0.90)=0.09.
\end{align*}

\begin{align*}
P(M_1\cap M_2^c)
&=P(M_1\cap M_2^c\mid L)P(L)\\
&\quad+P(M_1\cap M_2^c\mid L^c)P(L^c)\\
&=0.09(0.10)+0.09(0.90)\\
&=0.09.
\end{align*}

By Bayes' theorem,
\begin{align*}
P(L\mid M_1\cap M_2^c)
&=\frac{P(M_1\cap M_2^c\mid L)P(L)}
{P(M_1\cap M_2^c)}\\
&=\frac{0.09(0.10)}{0.09}\\
&=0.10.
\end{align*}

\[
\therefore\quad
P(L\mid M_1\cap M_2^c)=P(L)=0.10.
\]


\bigskip
\noindent
\textbf{(c)}

By the law of total probability,
\begin{align*}
P(M_1)
&=0.90(0.10)+0.10(0.90)\\
&=0.18,\\[4pt]
P(M_2)
&=0.90(0.10)+0.10(0.90)\\
&=0.18.
\end{align*}

\[
\therefore\quad
P(M_1)P(M_2)=0.18^2=0.0324.
\]

From part (b),
\[
P(M_1\cap M_2)=0.09.
\]

\[
\because\quad
P(M_1\cap M_2)=0.09
\neq 0.0324=P(M_1)P(M_2),
\]

\[
\therefore\quad
M_1\text{ and }M_2\text{ are not independent.}
\]

However, conditional independence gives
\[
P(M_1\cap M_2\mid L)
=P(M_1\mid L)P(M_2\mid L)
=0.81,
\]

\[
P(M_1\cap M_2\mid L^c)
=P(M_1\mid L^c)P(M_2\mid L^c)
=0.01.
\]

Moreover,
\begin{align*}
P(M_2\mid M_1)
&=\frac{P(M_1\cap M_2)}{P(M_1)}\\
&=\frac{0.09}{0.18}=0.50
\neq P(M_2)=0.18.
\end{align*}

But
\[
P(M_2\mid M_1,L)=P(M_2\mid L)=0.90,
\]

\[
P(M_2\mid M_1,L^c)=P(M_2\mid L^c)=0.10.
\]

Thus, the programs are conditionally independent
given the true email status, but not independent
without conditioning, because both outputs depend
on the same underlying email status.



\section*{Problem4}

Consider a simplified model of an AI evaluation. In one run,
an AI system answers two questions. Write $C$ for correct
and $I$ for incorrect; the ordered outcomes $CC,CI,IC,II$
also denote the corresponding events.

In parts (a) and (b), each answer is correct with probability
$1/2$, independently of the other, so the four outcomes
are equally likely.

An evaluator checks both answers against an answer key
and may publish one of the following reports.
Let $M_C$ and $M_I$ denote the events that the
corresponding reports are published:
\[
\begin{aligned}
M_C &: \quad \text{``At least one answer is correct''},\\
M_I &: \quad \text{``At least one answer is incorrect''}.
\end{aligned}
\]

At most one report is published per run, and every
published report must be true.

\textbf{(a)}
Initially, every run is reported. For $CC$, the evaluator
publishes $M_C$; for $II$, the evaluator publishes $M_I$.
For either $CI$ or $IC$, a fair coin is used to choose
between the two reports, each with probability $1/2$.

Using Bayes' rule, find
\[
P(CC\mid M_C)
\qquad\text{and}\qquad
P(II\mid M_I).
\]

\textbf{(b)}
Replace the rule in part (a) by any randomized reporting
rule based on the outcome. The evaluator may choose not
to publish a report for any of the four outcomes.
The two reports must each have positive probability
and must satisfy
\[
P(CC\mid M_C)\geq\frac23,
\qquad
P(II\mid M_I)\geq\frac34.
\]

What is the smallest possible probability that no report
is published? Specify a reporting rule for each of
$CC,CI,IC,II$ that attains your answer, and prove that
no admissible rule can do better.

\textbf{(c)}
Remove the independence assumption, but keep each
answer's probability of being correct equal to $1/2$.
The evaluator may choose a new randomized reporting
rule, but must publish exactly one truthful report
per run, with both report types having positive probability.

Let $t=P(CC)$. Determine all values of $t$ for which
such a rule satisfies both posterior inequalities
in part (b).

Prove necessity and sufficiency, giving a valid rule
for every claimed value of $t$. The rule may depend on $t$.

\section*{Solution}

\subsection*{(a)}

\[
P(CC)=P(CI)=P(IC)=P(II)=\frac14.
\]

\[
\begin{aligned}
P(M_C\mid CC)&=1,\\
P(M_C\mid CI)&=\frac12,\\
P(M_C\mid IC)&=\frac12,\\
P(M_C\mid II)&=0.
\end{aligned}
\]

By the law of total probability,
\[
\begin{aligned}
P(M_C)
&=\sum_{\omega\in\{CC,CI,IC,II\}}
P(M_C\mid\omega)P(\omega)\\
&=1\cdot\frac14
+\frac12\cdot\frac14
+\frac12\cdot\frac14
+0\cdot\frac14\\
&=\frac12.
\end{aligned}
\]

By Bayes' rule,
\[
\begin{aligned}
P(CC\mid M_C)
&=\frac{P(M_C\mid CC)P(CC)}{P(M_C)}\\
&=\frac{1\cdot\frac14}{\frac12}\\
&=\frac12.
\end{aligned}
\]

Similarly,
\[
\begin{aligned}
P(M_I)
&=0\cdot\frac14
+\frac12\cdot\frac14
+\frac12\cdot\frac14
+1\cdot\frac14\\
&=\frac12.
\end{aligned}
\]

\[
\begin{aligned}
P(II\mid M_I)
&=\frac{P(M_I\mid II)P(II)}{P(M_I)}\\
&=\frac{1\cdot\frac14}{\frac12}\\
&=\frac12.
\end{aligned}
\]

\[
\therefore\quad
P(CC\mid M_C)=P(II\mid M_I)=\frac12.
\]


\subsection*{(b)}

Let $N$ denote the event that no report is published.

\[
P(CC\mid M_C)\geq\frac23
\]

\[
\implies
\frac{P(CC\cap M_C)}{P(M_C)}
\geq\frac23
\]

\[
\implies
P(M_C)\leq\frac32P(CC\cap M_C).
\]

\[
\because\quad
P(CC\cap M_C)\leq P(CC)=\frac14
\]

\[
\therefore\quad
P(M_C)\leq\frac32\cdot\frac14=\frac38.
\]

Similarly,
\[
P(II\mid M_I)\geq\frac34
\]

\[
\implies
P(M_I)\leq\frac43P(II\cap M_I).
\]

\[
\because\quad
P(II\cap M_I)\leq P(II)=\frac14
\]

\[
\therefore\quad
P(M_I)\leq\frac43\cdot\frac14=\frac13.
\]

\[
\because\quad M_C\cap M_I=\varnothing
\]

\[
\begin{aligned}
P(N)
&=1-P(M_C)-P(M_I)\\
&\geq1-\frac38-\frac13\\
&=\frac7{24}.
\end{aligned}
\]

\textbf{Achievability:}

Consider the following reporting rule:
\[
\renewcommand{\arraystretch}{1.5}
\begin{array}{c|ccc}
\omega
&P(M_C\mid\omega)
&P(M_I\mid\omega)
&P(N\mid\omega)\\
\hline
CC&1&0&0\\
CI&\frac14&\frac16&\frac7{12}\\
IC&\frac14&\frac16&\frac7{12}\\
II&0&1&0
\end{array}
\]

\[
\begin{aligned}
P(M_C)
&=\frac14\left(1+\frac14+\frac14\right)\\
&=\frac38,
\end{aligned}
\]

\[
\begin{aligned}
P(M_I)
&=\frac14\left(\frac16+\frac16+1\right)\\
&=\frac13.
\end{aligned}
\]

\[
P(CC\mid M_C)
=\frac{P(CC\cap M_C)}{P(M_C)}
=\frac{\frac14}{\frac38}
=\frac23.
\]

\[
P(II\mid M_I)
=\frac{P(II\cap M_I)}{P(M_I)}
=\frac{\frac14}{\frac13}
=\frac34.
\]

\[
\begin{aligned}
P(N)
&=1-P(M_C)-P(M_I)\\
&=1-\frac38-\frac13\\
&=\frac7{24}.
\end{aligned}
\]

\[
\therefore\quad
\min P(N)=\frac7{24}.
\]


\subsection*{(c)}

Let
\[
t=P(CC).
\]

\[
P(C_1)=P(C_2)=\frac12,
\]
where $C_1,C_2$ denote correctness of the first
and second answers.

\[
\begin{aligned}
P(CI)&=P(C_1)-P(CC)=\frac12-t,\\
P(IC)&=P(C_2)-P(CC)=\frac12-t.
\end{aligned}
\]

\[
\begin{aligned}
P(II)
&=1-P(CC)-P(CI)-P(IC)\\
&=1-t-2\left(\frac12-t\right)\\
&=t.
\end{aligned}
\]

Thus,
\[
\begin{aligned}
P(CC)&=t,\\
P(CI)&=\frac12-t,\\
P(IC)&=\frac12-t,\\
P(II)&=t,
\end{aligned}
\qquad
0\leq t\leq\frac12.
\]

\textbf{Necessity:}

Let
\[
p=P(M_C).
\]

Exactly one truthful report is published per run:
\[
P(M_I)=1-p,
\qquad 0<p<1.
\]

\[
P(M_C\mid CC)=1,
\qquad
P(M_I\mid II)=1.
\]

\[
\implies
P(CC\cap M_C)=t,
\qquad
P(II\cap M_I)=t.
\]

The first posterior constraint gives
\[
P(CC\mid M_C)=\frac{t}{p}\geq\frac23
\]

\[
\implies\quad
p\leq\frac32t.
\]

The second posterior constraint gives
\[
P(II\mid M_I)=\frac{t}{1-p}\geq\frac34
\]

\[
\implies\quad
1-p\leq\frac43t
\]

\[
\implies\quad
p\geq1-\frac43t.
\]

\[
\therefore\quad
1-\frac43t\leq p\leq\frac32t.
\]

\[
\implies
1-\frac43t\leq\frac32t
\]

\[
\implies
1\leq\frac{17}{6}t
\]

\[
\implies
t\geq\frac6{17}.
\]

\[
\because\quad t\leq\frac12
\]

\[
\therefore\quad
\frac6{17}\leq t\leq\frac12.
\]


\textbf{Sufficiency:}

For every
\[
t\in\left[\frac6{17},\frac12\right],
\]
define
\[
q(t)=
\begin{cases}
\dfrac{t}{2(1-2t)},
&\dfrac6{17}\leq t\leq\dfrac25,\\[8pt]
1,
&\dfrac25<t\leq\dfrac12.
\end{cases}
\]

Then
\[
0\leq q(t)\leq1.
\]

Consider the reporting rule:
\[
\renewcommand{\arraystretch}{1.5}
\begin{array}{c|cc}
\omega
&P(M_C\mid\omega)
&P(M_I\mid\omega)\\
\hline
CC&1&0\\
CI&q(t)&1-q(t)\\
IC&q(t)&1-q(t)\\
II&0&1
\end{array}
\]

Each row sums to $1$, and every report is truthful.

\textbf{Case 1:}
\[
\frac6{17}\leq t\leq\frac25.
\]

\[
\begin{aligned}
P(M_C)
&=t+2\left(\frac12-t\right)q(t)\\
&=t+(1-2t)\frac{t}{2(1-2t)}\\
&=\frac32t.
\end{aligned}
\]

\[
P(M_I)=1-\frac32t.
\]

\[
P(CC\mid M_C)
=\frac{t}{\frac32t}
=\frac23.
\]

\[
P(II\mid M_I)
=\frac{t}{1-\frac32t}.
\]

\[
\because\quad
t\geq\frac6{17}
\]

\[
\implies
\frac{17}{2}t\geq3
\]

\[
\implies
4t\geq3-\frac92t
\]

\[
\implies
\frac{t}{1-\frac32t}\geq\frac34.
\]

\[
\therefore\quad
P(CC\mid M_C)\geq\frac23,
\qquad
P(II\mid M_I)\geq\frac34.
\]


\textbf{Case 2:}
\[
\frac25<t\leq\frac12.
\]

\[
q(t)=1.
\]

\[
\begin{aligned}
P(M_C)
&=P(CC)+P(CI)+P(IC)\\
&=t+2\left(\frac12-t\right)\\
&=1-t.
\end{aligned}
\]

\[
P(M_I)=P(II)=t.
\]

\[
P(CC\mid M_C)=\frac{t}{1-t}.
\]

\[
\because\quad t>\frac25
\]

\[
\implies
5t>2
\]

\[
\implies
3t>2(1-t)
\]

\[
\implies
\frac{t}{1-t}>\frac23.
\]

Furthermore,
\[
P(II\mid M_I)=\frac{t}{t}=1\geq\frac34.
\]

\[
\therefore\quad
P(CC\mid M_C)\geq\frac23,
\qquad
P(II\mid M_I)\geq\frac34.
\]

In both cases,
\[
P(M_C)>0,
\qquad
P(M_I)>0.
\]

Combining necessity and sufficiency,
\[
\therefore\quad
t\in\left[\frac6{17},\frac12\right].
\]




\end{document}
